
\scanpage{146}
Question : \unsure{Sont uniformément convexes|Sont-ce des convexes|Sont-ils convexes|Dans quel cas (uniforme convexité} \unsure{les|ces|des} \unsure{espaces|esp.} $L^p(E,F)$ ? pour $0<p<\infty$ ? Pas d'importance.

2) Pour la continuité de $h\to h^r$, on est ramené au cas $0\leq r<1$ (car c'est trivial pour $r$ entier) Il semble possible que

\[
\bigl\|\lvert {\unsure{h|H}}^r-k^r\rvert\bigr\|^{1/r}\leq\|h-k\|\qquad\text{pour }0\leq r<1.
\]

Il \unsure{revient au $=$ de prouver|revient au $=$ de présumer} que

\[
S_{p/r}(h^r-k^r)=S_p\bigl(\lvert h^r-k^r\rvert^{1/r}\bigr)
\]

est fonction croissante de $r$, ($h$ et $k$ hermitiens positifs, $p$ quelconque).

On peut faire aussi $p=1$.
% archive sheet number: 147
